% 方法章节管线：定义研究合同，解释三个贡献及其代码落点，最后给出尚待执行的消融方案。
% 公式是对已有实现的统一描述，不宣称新增算法、因果识别或已证明的性能提升。
\section{方法创新：机制、异质性与证据的递归耦合}
\label{sec:innovation}
\subsection{贡献一：以研究合同统一语义与计算}
AURORA 将一个研究对象概括为
\begin{equation}
\mathcal{H}=(m,\mathcal{F},h,\Omega,\mathcal{A},\ell),
\end{equation}
其中 $m$ 为机制描述，$\mathcal{F}$ 为操作化表达式集合，$h$ 为冻结预测周期，$\Omega$ 为覆盖条件，$\mathcal{A}$ 为必要断言，$\ell$ 为父子谱系。测量时再绑定配置、数据和代码身份。该符号是对现有 \code{Structure} 与冻结产物的归纳，不新增另一套运行协议。

设计价值在于\important{让语义假设与统计检验保持对应}：提案器需要说明机制，操作化角色需要给出合法表达式，审查角色需要检查断言是否可测。某个表达式 IC 较高，并不能跳过该机制必要条件的核对。10 类机制与 7 类形式提供可导航的初始目录，递归队列继续派生新的研究对象。

\begin{center}
\setlength{\fboxsep}{9pt}
\fcolorbox{aurorablue!25}{codegray}{\begin{minipage}{.92\textwidth}
\textcolor{aurorablue}{\textbf{接口约束即方法的一部分}}\\[0.3em]
机制描述 $\rightarrow$ 受限 DSL 与断言校验 $\rightarrow$ 实验冻结 $\rightarrow$ 确定性测量。\\
模型组织研究问题；计算模块产生数值与检验状态；审计产物保留两者的对应关系。
\end{minipage}}
\end{center}

\subsection{贡献二：将局部异质性转换为有谱系的新假设}
整体效应弱时，可能仍存在值得检验的状态依赖。系统使用训练段线索构造条件或方向子代。为说明这种变换，可将条件子代概括为
\begin{equation}
 f_{\mathrm{child}}(x)=d\,f_{\mathrm{parent}}(x)\,G(x),\qquad
 G(x)=\prod_{j=1}^{k}\mathbf{1}\{x_j\in I_j\},\quad d\in\{-1,+1\}.
\end{equation}
这里 $I_j$ 为训练选择并冻结的条件区域，$d$ 为训练确定的方向。公式表示一种条件化方式；实际实现还保留覆盖掩码，并支持其他演化操作，不能据此把覆盖外的未知收益填成零。

\important{完整路径比单个阈值更适合作为研究对象。}多项条件共同定义局部状态，子代携带父结构、来源条件和操作记录，重新接受审查与测量。其解释从“一个信号是否有效”细化为“这个信号在什么状态、什么方向和什么周期下值得继续检验”。森林所选条件仍是探索假设，不因模型或树结构给出而成为因果结论。

\clearpage
\subsection{贡献三：将证据冲突保留为下一轮研究输入}
证据记录同时承担评估与计划两个职责。前者保留结果状态和必要断言，后者组织可继续探索的线索；二者通过明确接口连接。\code{engine/memory.py} 的后验服务研究计划，\code{engine/confirmation.py} 负责另行检查正式确认资格。模型不能通过改写解释来修改已测数值或确认状态。

\begin{table}[ht]
\centering\small
\begin{tabularx}{\textwidth}{@{}p{.20\textwidth}XX@{}}
\toprule
\textbf{观察} & \textbf{可采取的研究动作} & \textbf{必须保留的证据边界}\\
\midrule
整体较弱，局部有线索 & 登记条件子代，再次冻结与测量 & 条件由训练选择，不继承父代确认资格\\
训练呈反方向 & 登记方向诊断，满足队列约束时提出反向候选 & 后续段方向冲突必须保留\\
已有支持也有冲突 & 在记录中同时呈现，区分观察、剔除与后续候选 & 发现后解释不能回写为事前先验\\
检验无效或未完成 & 记录未知与阻塞原因，按流程补充测量 & 未测试不能转写为支持证据\\
\bottomrule
\end{tabularx}
\caption{证据反馈的是研究动作；确认资格由独立流程判定。动作受预算、深度和准入条件限制。}
\end{table}

上述设计形成三个可审查的不变量：\textbf{选择来源有记录、研究对象可冻结、确认状态有独立入口}。角色白名单限制输入，冻结指纹检查一致性，检查点保留已完成测量。它们共同支撑递归过程的可追溯性，但不等同于完整信息安全认证或对所有统计风险的自动消除。

\subsection{贡献与实现证据的对应关系}
\begin{tabularx}{\textwidth}{@{}p{.18\textwidth}XX@{}}
\toprule
\textbf{贡献} & \textbf{主要代码落点} & \textbf{现有证据与下一步验证}\\
\midrule
研究合同 & \code{catalog.py}、\code{dsl.py}、\code{llm.py}、\code{pipeline.py} & 已有冻结表达式和断言产物；待比较机制约束对合同有效率的影响\\
异质性演化 & \code{discovery.py}、\code{evolution.py}、\code{cycle.py} & 归档含 6 个子代，截图含条件案例；待执行同预算演化对照\\
证据记忆 & \code{memory.py}、\code{laws.py}、\code{confirmation.py} & 已有结果及规律记录；待量化记忆对重复探索和复核成本的影响\\
\bottomrule
\end{tabularx}

\vspace{0.5em}
这三项贡献强调模块之间的协议与反馈关系。当前材料支持实现与案例层面的论证，不主张各组成算法均由本项目首创，也不以训练 IC 证明整个系统优越性。

\clearpage
\subsection{如何检验创新是否带来增益：同预算消融方案}
\textcolor{aurorablue}{\textbf{本小节为后续实验设计，尚未执行，不计入当前结果。}}为了区分工程闭环与方法增益，建议在同一数据版本、计算环境、测量口径和研究预算下比较以下配置。

\begin{table}[ht]
\centering\small
\begin{tabularx}{\textwidth}{@{}p{.17\textwidth}XX@{}}
\toprule
\textbf{对照配置} & \textbf{保留与移除} & \textbf{拟回答的问题}\\
\midrule
A：初始探索 & 保留机制合同、审查和测量；关闭自适应后续队列 & 固定初始坐标能提供怎样的研究基线？\\
B：加入演化 & 相同合同与检验，开放方向、条件及交互等子代；关闭研究记忆反馈 & 演化是否增加可复核的新假设及后续稳定性？\\
C：完整闭环 & 在 B 上加入研究记忆与探索调度 & 记忆是否减少重复探索、改善单位预算产出？\\
\bottomrule
\end{tabularx}
\caption{拟议递增消融。实现配置时需记录实际开关，不能假设现有 CLI 已包含全部消融选项。}
\end{table}

\textbf{控制条件。}冻结相同的数据资格规则、时间分段、模型版本与调用设置；为各组设置相同的候选尝试上限、模型调用上限和计算时间上限，先触及任一上限即停止。各配置可能因初始队列耗尽而提前结束，因此同时披露实际资源消耗和未使用预算，不补造初始候选以填满预算。使用预先登记的多次随机种子复跑；分别报告拒绝、回退、完成和继承记录。探索结果可以用于诊断，独立确认样本不能反复用于调参。

\textbf{评价指标。}优先使用合同有效率、完成测量的单位成本、唯一冻结假设数量、父子变化可解释率和人工复核时间。若进行独立确认，需事前冻结全部比较配置及待检候选，统一约定确认检验族与多重比较规则，避免跨组反复使用保留样本。确认后再报告满足前提的通过数量与区间；不以最高训练 IC 作为唯一评价标准。

\begin{center}
\setlength{\fboxsep}{9pt}
\fcolorbox{auroragold!45}{auroragold!6}{\begin{minipage}{.92\textwidth}
\textcolor{aurorablue}{\textbf{可被否定的增益主张}}\\[0.3em]
若加入演化或记忆后，研究重复度、复核成本或独立确认产出未改善，应报告无增益，并检查队列策略及记忆可识别性。这样才能将“具备自演化能力”推进为“自演化具有可测量价值”。
\end{minipage}}
\end{center}

\subsection{从研究原型到团队工具}
当前适合以本地工作台和 SDK 接入研究团队的授权数据，先交付一条可复核的研究链。试点评估数据接入时间、研究记录完整率、异常恢复与人工审查负担。权限、团队协作和运维能力按试点需求补齐；商业化指标与投资有效性分别验证。
\clearpage